\chapter{Literature Survey}
\label{chap:literature}

\section{Evolution of RF Sensing: RSSI vs. CSI}
Device-free indoor sensing originated with the analysis of Received Signal Strength Indicator (RSSI) values extracted from standard 802.11 frames \cite{patwari2014breathing, lara2013survey}. Patwari \textit{et al.} demonstrated that an individual traversing an RF link induces shadowing effects that modulate RSSI. However, RSSI is fundamentally a single macroscopic scalar value per packet, aggregating all multipath reflections constructively and destructively across the entire frequency band. Consequently, RSSI exhibits severe multipath fading, high temporal variance, and vulnerability to background ambient interference, rendering it incapable of resolving millimeter-scale chest wall breathing motions \cite{xiao2013pilot, zeng2019farsense}.

The release of open-source CSI tools for the Intel 5300 NIC by Halperin \textit{et al.} \cite{halperin2011tool} revolutionized wireless sensing. By tapping into physical-layer baseband receivers, researchers gained access to subcarrier-resolved Channel Frequency Response (CFR) matrices. Unlike RSSI, Channel State Information provides fine-grained amplitude and phase information across orthogonal subcarriers, revealing sub-wavelength spatial perturbations caused by human motion.

\section{Subcarrier-Resolved Physical-Layer Sensing Systems}
With the availability of CSI, numerous activity recognition frameworks emerged. CARM \cite{wang2015understanding} and RT-Fall \cite{wang2017rt} quantified the relationship between human kinematic velocity and CSI Doppler frequency shifts, demonstrating device-free fall detection and activity classification. Xiao \textit{et al.} introduced Pilot \cite{xiao2013pilot}, utilizing CSI amplitude correlation for passive localization. However, these systems depended entirely on specialized mini-PCIe NICs, modified Linux kernel drivers, and dedicated desktop workstations.

\section{Contactless Respiration and Vital Sign Monitoring}
Contactless vital sign tracking leverages the physiological reality that thoracic displacement during normal tidal breathing (1--5~mm) alters the path length of reflected RF rays \cite{zeng2019farsense, liu2015tracking}. Phase-based CSI sensing offers superior sensitivity over amplitude, as subcarrier phase shift $\Delta \phi$ is inversely proportional to carrier wavelength $\lambda$:
\begin{equation}
\Delta \phi = 2\pi \frac{\Delta d}{\lambda}
\label{eq:phase_path_lit}
\end{equation}
where $\lambda \approx 12.5$~cm for 2.4~GHz Wi-Fi. Benchmark systems such as PhaseBeat \cite{liu2015tracking} and FarSense \cite{zeng2019farsense} utilized CSI phase difference and ratio across multiple receiver antennas to cancel carrier frequency offsets (CFO). However, these platforms required multiple external directional antennas, dual-band transceivers, and complex offline MATLAB processing.

\section{Edge-Scale IoT Sensing Architectures}
Recently, low-cost microcontrollers, notably the Espressif ESP32, have introduced firmware-level physical-layer CSI extraction \cite{hernandez2020lightweight}. Unlike PCIe NICs, the ESP32 is a self-contained, low-power SoC operating at 240~MHz with integrated Wi-Fi. Early exploratory investigations validated the feasibility of capturing CSI on ESP32 \cite{hernandez2020lightweight, guo2022device}. However, existing ESP32 prototypes either dump raw packets to external cloud servers for heavy post-processing or suffer from threshold instability under ambient RF drift.

\section{Comparative Analysis and Research Gap}
Table~\ref{tab:sota_comparison} presents a comprehensive comparative analysis of existing Wi-Fi sensing systems against the proposed framework.

\begin{table}[ht]
\centering
\caption{Comparison of Proposed Sensing Framework Against State-of-the-Art Systems}
\label{tab:sota_comparison}
\begin{tabular}{|l|c|c|c|c|c|}
\hline
\textbf{System} & \textbf{Hardware} & \textbf{Subcarriers} & \textbf{Cost (USD)} & \textbf{Vital Sensing} & \textbf{Edge Capable} \\
\hline
Pilot \cite{xiao2013pilot} & Atheros AR9380 & 56 & $\approx \$120$ & No & Host PC \\
CARM \cite{wang2015understanding} & Intel 5300 & 30 & $\approx \$180$ & No & Linux PC \\
PhaseBeat \cite{liu2015tracking} & Intel 5300 & 30 & $\approx \$180$ & Respiration & Linux PC \\
FarSense \cite{zeng2019farsense} & Intel 5300 (3 Ant.) & 30 & $\approx \$220$ & Respiration & Server/MATLAB \\
CSI-Net \cite{wang2019csinet} & Intel 5300 + GPU & 30 & $\approx \$1500$ & No & GPU Server \\
\hline
\textbf{Proposed} & \textbf{ESP32 SoC} & \textbf{52 (I/Q)} & \textbf{$\approx \$4$} & \textbf{Yes (BPM)} & \textbf{Edge Real-Time} \\
\hline
\end{tabular}
\end{table}

\noindent \textbf{Research Gap Identified:} Existing systems require expensive, bulky hardware, multi-antenna configurations, or cloud-bound deep neural networks. There is a decisive gap for a unified, edge-native mathematical architecture capable of achieving high-accuracy presence detection and respiration monitoring on a single-antenna \$4 microcontroller in real time.
